\documentclass[11pt]{article}

\usepackage[T1]{fontenc}
\usepackage[utf8]{inputenc}
\usepackage[margin=1in]{geometry}
\usepackage{amsmath,amssymb,amsthm}
\usepackage[shortlabels]{enumitem}
\usepackage{microtype}

\setlength{\parindent}{0pt}
\setlength{\parskip}{0.65em}
\setlist[enumerate]{leftmargin=*,itemsep=0.45em,topsep=0.35em}

\newcommand{\N}{\mathbb{N}}
\newcommand{\Z}{\mathbb{Z}}
\newcommand{\Q}{\mathbb{Q}}
\newcommand{\R}{\mathbb{R}}
\newcommand{\problem}[2]{\subsection*{Problem #1 \hfill \normalfont[#2 points]}}

\begin{document}
\pagestyle{empty}

\begin{center}
    {\Large\bfseries 21-128 and 15-151}\\[3pt]
    {\LARGE\bfseries Problem Sheet 1}\\[7pt]

Solutions to the following six exercises and optional bonus problem are to be 
submitted to Gradescope by

\textbf{11 PM on Wednesday, the 2nd of September, 2026.}
\end{center}

\hrule
\vspace{0.6em}

\textbf{Instructions.}
This assignment has 30 points, together with an optional 1-point bonus problem.
Unless a problem explicitly says otherwise, every answer must be justified by a proof.
A counterexample must satisfy all of the hypotheses and must be checked explicitly.
You may cite results proved in lecture or recitation, provided that you identify the 
result you are using.


\problem{1}{4}
Let $a,b\in\Z_{>0}$.
\begin{enumerate}[(a)]
    \item Suppose that $a\mid b$. Prove that
    \[
        a\mid a^2+b^2.
    \]

    \item Is the converse true? In other words, if
    \[
        a\mid a^2+b^2,
    \]
    must $a\mid b$? Prove your answer.
\end{enumerate}


\problem{2}{4}
Find, with proof, all real solutions of
\[
    5=2\sqrt{x}+\sqrt{10-x}.
\]
Your solution should identify the domain of the equation, explain whether each squaring step is reversible, and check every candidate in the original equation.


\problem{3}{5}
Let $D(m,n)$ be the predicate ``$m$ divides $n$,'' where all variables range over $\Z_{>0}$.
Determine whether each proposition below is true or false.
\[
\begin{array}{ll}
\text{(a)} & \forall m\,\exists n\; D(m,n),\\[2pt]
\text{(b)} & \exists n\,\forall m\; D(m,n),\\[2pt]
\text{(c)} & \forall n\,\exists m\; D(m,n),\\[2pt]
\text{(d)} & \exists m\,\forall n\; D(m,n).
\end{array}
\]
For every true proposition, give a valid witness or a rule for choosing a witness.
For every false proposition, state and prove its negation.


\problem{4}{5}
\begin{enumerate}[(a)]
    \item Prove that
    \[
        (p\Rightarrow q)\vee(p\Rightarrow r)
        \qquad\text{and}\qquad
        p\Rightarrow(q\vee r)
    \]
    are logically equivalent. You may use a truth table or standard logical equivalences, but every step must be clear.

    \item Let $P$ and $Q$ be predicates on a set $S$. Prove that
    \[
        \bigl(\forall x\in S\;P(x)\bigr)
        \vee
        \bigl(\forall x\in S\;Q(x)\bigr)
        \quad\Longrightarrow\quad
        \forall x\in S\;\bigl(P(x)\vee Q(x)\bigr).
    \]
    Then give explicit choices of $S$, $P$, and $Q$ showing that the converse implication is false.
\end{enumerate}


\problem{5}{7}
Consider the equation
\[
    x^4y+ay+x=0,
\]
where $a,x,y\in\R$.
\begin{enumerate}[(a)]
    \item Show that the proposition
    \[
        (\forall a\in\R)(\forall x\in\R)(\exists!\,y\in\R)
        \bigl(x^4y+ay+x=0\bigr)
    \]
    is false. Give one pair $(a,x)$ for which no real solution $y$ exists and another pair $(a,x)$ for which more than one real solution $y$ exists.

    \item Determine, with proof, the set of all real numbers $a$ for which
    \[
        (\forall x\in\R)(\exists!\,y\in\R)
        \bigl(x^4y+ay+x=0\bigr)
    \]
    is true.
\end{enumerate}


\problem{6}{5}
Let $r\in\Q\setminus\{0\}$ and let $a,b\in\R\setminus\Q$.
Determine which of the following expressions are irrational for \emph{every} permitted choice of $r,a,b$:
\[
    a+r,
    \qquad
    ar,
    \qquad
    a+b,
    \qquad
    ab.
\]
For each expression that must be irrational, give a proof.
For each expression that need not be irrational, give an explicit counterexample and verify that all numbers in the example satisfy the hypotheses.


\newpage
\subsection*{Bonus Problem \hfill \normalfont[1 bonus point]}
Katie, Hasita, and Michelle are perfect logicians. Each has a positive integer written on their hat. 
The numbers are not necessarily distinct, and it is common knowledge that two of the three numbers 
add to the third. Each logician can see the other two hats but not their own.

The following conversation is public:
\begin{quote}
\textbf{Katie:} I do not know my number.\\
\textbf{Hasita:} I do not know my number.\\
\textbf{Michelle:} I do not know my number.\\
\textbf{Katie:} Now I know my number. It is $50$.
\end{quote}

\begin{enumerate}[(a)]
    \item Determine the numbers on Hasita's and Michelle's hats, and justify every inference.

    \item At Michelle's first turn---after they have heard Katie and Hasita each say ``I do not know''---characterize all ordered pairs $(u,v)$ that they could see on Katie's and Hasita's hats for which they would be able to determine their own number. Your condition should be both necessary and sufficient.
\end{enumerate}

\end{document}
